\section{Lan truyền ngược theo thời gian}
\label{sec:bptt}

Mạng hồi quy (RNN) đưa vấn đề triệt tiêu gradient sang một dạng nghiêm trọng hơn
hẳn. Lý do không nằm ở chiều sâu lớn hơn, mà ở một tính chất cấu trúc sẽ được chỉ
rõ ở Mục~V.3.

\begin{notebox}
\textbf{Quy ước ký hiệu cho mục này.} $T$ (chữ nghiêng, in thường trong công
thức) luôn là \emph{số bước thời gian}. Phép chuyển vị luôn viết bằng $\top$, ví
dụ $\mathbf{W}_{hh}^{\top}$. Trong toàn mục này không bao giờ dùng $W^{T}$ để chỉ
chuyển vị --- ký hiệu đó sẽ mang nghĩa ``luỹ thừa $T$''.
\end{notebox}

\subsection{Mô hình và ký hiệu}

Xét RNN với trạng thái ẩn $\mathbf{h}_t\in\mathbb{R}^{n}$:
\begin{equation}
  \mathbf{u}_t = \mathbf{W}_{hh}\mathbf{h}_{t-1}
               + \mathbf{W}_{xh}\mathbf{x}_t + \mathbf{b}_h,
  \qquad
  \mathbf{h}_t = \tanh\!\left(\mathbf{u}_t\right).
  \label{eq:rnn}
\end{equation}
Ta tách tiền kích hoạt $\mathbf{u}_t$ ra khỏi trạng thái $\mathbf{h}_t$ để giữ
đúng quy ước của Mục~\ref{sec:cosotoan}, nơi $\boldsymbol{\delta}$ luôn được định
nghĩa theo \emph{tiền} kích hoạt. Đặt
\[
  \boldsymbol{\delta}_t := \frac{\partial\mathcal{L}_T}{\partial\mathbf{u}_t},
  \qquad
  \mathbf{g}_t := \frac{\partial\mathcal{L}_T}{\partial\mathbf{h}_t},
\]
với $\mathcal{L}_T$ là mất mát tại bước cuối $T$. Vì $\tanh$ tác động theo từng
phần tử và $\tanh' = 1 - \tanh^2$, ta có sẵn
$\tanh'(\mathbf{u}_t) = 1 - \mathbf{h}_t^2$ (bình phương theo từng phần tử) ---
một tiện lợi tính toán: đạo hàm biểu diễn được qua chính đầu ra.

\subsection{Công thức đệ quy}

\begin{infobox}{Đệ quy BPTT theo tiền kích hoạt}
Với quy ước ở Mục~\ref{sec:cosotoan},
\begin{equation}
  \boldsymbol{\delta}_t
  = \operatorname{diag}\!\left(1-\mathbf{h}_t^2\right)
    \mathbf{W}_{hh}^{\top}\,\boldsymbol{\delta}_{t+1},
  \label{eq:bpttrec}
\end{equation}
và do đó
\begin{equation}
  \boldsymbol{\delta}_1
  = \left(\prod_{t=2}^{T}
      \operatorname{diag}\!\left(1-\mathbf{h}_{t-1}^2\right)\mathbf{W}_{hh}^{\top}
    \right)\boldsymbol{\delta}_T .
  \label{eq:bpttprod}
\end{equation}
\end{infobox}

\noindent\textit{Chứng minh.}\;
Từ~\eqref{eq:rnn}, $\mathbf{u}_{t+1}$ phụ thuộc $\mathbf{u}_t$ qua
$\mathbf{u}_{t+1} = \mathbf{W}_{hh}\tanh(\mathbf{u}_t) + \cdots$, nên
\[
  \frac{\partial u_{t+1,m}}{\partial u_{t,i}}
  = W_{hh,\,mi}\left(1 - h_{t,i}^2\right).
\]
Quy tắc dây chuyền cho
\[
  \delta_{t,i}
  = \sum_m \delta_{t+1,m}\,W_{hh,\,mi}\left(1-h_{t,i}^2\right)
  = \left(1-h_{t,i}^2\right)\left[\mathbf{W}_{hh}^{\top}\boldsymbol{\delta}_{t+1}\right]_i ,
\]
đúng dạng~\eqref{eq:bpttrec}. Lặp lại từ $t=1$ tới $T$ cho~\eqref{eq:bpttprod}.
\hfill$\square$

\begin{notebox}
Nếu viết đệ quy theo \emph{trạng thái} $\mathbf{h}$ thay vì tiền kích hoạt, thứ
tự hai thừa số đảo lại:
$\mathbf{g}_t = \mathbf{W}_{hh}^{\top}\operatorname{diag}\!\left(1-\mathbf{h}_{t+1}^2\right)\mathbf{g}_{t+1}$,
vì khi đó phi tuyến nằm ở bước $t+1$ chứ không phải bước $t$. Hai cách viết hoàn
toàn tương đương và cho cùng một chặn; nhầm lẫn giữa chúng là lỗi phổ biến nhất
khi cài đặt BPTT bằng tay. Quy ước theo $\mathbf{u}$ ở~\eqref{eq:bpttrec} khớp
chính xác với dạng $\operatorname{diag}\cdot\mathbf{W}^{\top}$
của~\eqref{eq:deltarec}.
\end{notebox}

\subsection{Vì sao RNN tệ hơn MLP cùng độ sâu}

Áp dụng tính dưới nhân cho~\eqref{eq:bpttprod}, giống hệt cách làm
ở~\eqref{eq:governing}:
\begin{equation}
  \left\|\boldsymbol{\delta}_1\right\|
  \;\le\;
  \left(\sigma_{\max}\!\left(\mathbf{W}_{hh}\right)\cdot\max\tanh'\right)^{T-1}
  \left\|\boldsymbol{\delta}_T\right\|
  \;\le\;
  \sigma_{\max}\!\left(\mathbf{W}_{hh}\right)^{T-1}\left\|\boldsymbol{\delta}_T\right\| .
  \label{eq:bpttbound}
\end{equation}

Điểm cấu trúc then chốt nằm ở chỗ: trong~\eqref{eq:bpttprod}, \textbf{cùng một ma
trận $\mathbf{W}_{hh}$ được nhân lặp lại $T-1$ lần}. Đây là khác biệt cốt yếu so
với MLP.

Trong một MLP $L$ tầng, tích~\eqref{eq:central} gồm $L$ ma trận \emph{khác nhau}
$\mathbf{W}^{[1]},\dots,\mathbf{W}^{[L]}$. Các giá trị
$\sigma_{\max}(\mathbf{W}^{[k]})$ dao động quanh một giá trị trung bình, và tích
của chúng chịu một hiệu ứng trung bình hoá: một tầng có
$\sigma_{\max} = 0.9$ có thể được bù bởi một tầng có $\sigma_{\max} = 1.1$. Độ
lệch tích luỹ tăng theo $\sqrt{L}$ chứ không theo $L$.

Trong RNN, \emph{không có sự trung bình hoá nào cả}. Chặn
trong~\eqref{eq:bpttbound} là $\sigma_{\max}(\mathbf{W}_{hh})^{T-1}$ --- một cấp
số nhân thuần tuý với công bội cố định. Nếu
$\sigma_{\max}(\mathbf{W}_{hh}) = 0.9$, thì sau $T = 100$ bước thời gian hệ số là
$0.9^{99} \approx 2.95\times10^{-5}$, không có cơ chế nào kéo nó lại. Nói thẳng:

\begin{guidelinebox}{Kết luận so sánh}
Một RNN duỗi qua $T$ bước thời gian \textbf{tệ hơn ngặt} một MLP $T$ tầng, dù hai
mạng có cùng chiều sâu tính toán. Chia sẻ trọng số biến tích các ma trận độc lập
thành luỹ thừa của một ma trận duy nhất, loại bỏ hiệu ứng trung bình hoá và biến
mọi độ lệch nhỏ của $\sigma_{\max}$ khỏi $1$ thành phân kỳ hoặc triệt tiêu theo
cấp số nhân~\cite{BengioSimardFrasconi1994,Pascanu2013}.
\end{guidelinebox}

\subsection{Phân tích phổ: điều gì thực sự quyết định tốc độ suy giảm}

Chặn~\eqref{eq:bpttbound} dùng chuẩn phổ nên là một chặn \emph{trên}. Phân tích
riêng cho biết hành vi \emph{thực sự}.

Giả sử $\mathbf{W}_{hh}$ chéo hoá được:
$\mathbf{W}_{hh} = \mathbf{Q}\boldsymbol{\Lambda}\mathbf{Q}^{-1}$ với
$\boldsymbol{\Lambda} = \operatorname{diag}(\lambda_1,\dots,\lambda_n)$. Khi đó
\[
  \mathbf{W}_{hh}^{\top}
  = \left(\mathbf{Q}^{-1}\right)^{\!\top}\boldsymbol{\Lambda}\,\mathbf{Q}^{\top},
  \qquad\text{nên}\qquad
  \left(\mathbf{W}_{hh}^{\top}\right)^{T-1}
  = \left(\mathbf{Q}^{-1}\right)^{\!\top}\boldsymbol{\Lambda}^{\,T-1}\mathbf{Q}^{\top},
\]
vì các thừa số $\mathbf{Q}^{\top}(\mathbf{Q}^{-1})^{\top} = \mathbf{I}$ triệt tiêu
từng cặp ở giữa. Do $\boldsymbol{\Lambda}^{T-1} =
\operatorname{diag}(\lambda_1^{T-1},\dots,\lambda_n^{T-1})$, một thành phần
gradient nằm dọc hướng riêng thứ $i$ bị nhân đúng với $\lambda_i^{\,T-1}$:

\begin{infobox}{Suy giảm theo từng hướng riêng}
Thành phần gradient dọc hướng riêng ứng với $\lambda_i$ suy giảm theo cấp số nhân
$\left|\lambda_i\right|^{T-1}$. Mọi hướng có $|\lambda_i| < 1$ tiêu biến; chỉ các
hướng có $|\lambda_i|\ge1$ còn sống sót qua nhiều bước thời gian. Sau đủ nhiều
bước, gradient bị \emph{ép} vào không gian con sinh bởi các hướng riêng trội.
\end{infobox}

\subsubsection{Hai điều kiện cần nói rõ}

\textbf{Thứ nhất, chéo hoá được là một giả thiết.} Không phải mọi ma trận đều
chéo hoá được; với ma trận khuyết, phải dùng dạng chuẩn Jordan và các khối Jordan
sinh thêm thừa số đa thức theo $T$. Các thừa số đa thức này không đổi bản chất
cấp số nhân của kết luận, nhưng làm công thức không còn gọn.

\textbf{Thứ hai, và quan trọng hơn: với $\mathbf{W}_{hh}$ không chuẩn tắc}
(\emph{non-normal}, tức $\mathbf{W}_{hh}\mathbf{W}_{hh}^{\top}\neq
\mathbf{W}_{hh}^{\top}\mathbf{W}_{hh}$), bán kính phổ chỉ chi phối hành vi
\emph{tiệm cận} khi $T\to\infty$. Ở các $T$ hữu hạn, biên độ có thể \emph{tăng
vọt tạm thời} vượt xa mức mà $\max_i|\lambda_i|$ gợi ý, do các hướng riêng không
trực giao nhau. Một ma trận có mọi $|\lambda_i| < 1$ vẫn có thể khuếch đại
gradient hàng chục lần ở những bước trung gian trước khi cuối cùng suy giảm. Đây
là lý do $\sigma_{\max}$ --- chứ không phải $|\lambda_{\max}|$ --- mới là đại
lượng dùng để \emph{bảo đảm} an toàn.

Vì vậy phát biểu đúng của Pascanu và cộng sự~\cite{Pascanu2013} là bất đối xứng,
và cần được trích dẫn đúng như vậy:

\begin{guidelinebox}{Phát biểu chính xác --- chú ý tính bất đối xứng}
\begin{itemize}[leftmargin=*, itemsep=2pt]
  \item $\sigma_{\max}(\mathbf{W}_{hh}) < 1$ là điều kiện \textbf{đủ} để gradient
    triệt tiêu. Đây là một bảo đảm: nếu chuẩn phổ nhỏ hơn $1$, gradient
    \emph{chắc chắn} tiêu biến theo cấp số nhân.
  \item $\left|\lambda_{\max}(\mathbf{W}_{hh})\right| > 1$ là điều kiện
    \textbf{cần nhưng không đủ} để gradient bùng nổ. Bán kính phổ lớn hơn $1$ mở
    ra khả năng bùng nổ nhưng không bảo đảm nó xảy ra: hướng riêng trội có thể
    gần như trực giao với gradient thực tế, hoặc thừa số
    $\operatorname{diag}(1-\mathbf{h}_t^2)\le1$ có thể ghìm tích lại.
\end{itemize}
Hệ quả thực hành: dùng $\sigma_{\max}$ để chẩn đoán triệt tiêu, và đừng suy ra
``an toàn'' chỉ từ việc $|\lambda_{\max}| < 1$.
\end{guidelinebox}

\subsection{Quên thảm hoạ các phụ thuộc xa}

Hệ quả thực tiễn của~\eqref{eq:bpttbound} là hiện tượng \emph{quên thảm hoạ}
(catastrophic forgetting) các phụ thuộc dài hạn. Xét câu:

\begin{quote}
\itshape
``Tôi lớn lên ở Huế, \ldots{} (năm mươi từ xen giữa) \ldots{} nên tôi nói giọng
\underline{\phantom{Huế}}.''
\end{quote}

Để dự đoán từ cuối, mô hình phải dùng thông tin từ bước $t=5$ tại bước $t=60$.
Gradient của mất mát tại $t=60$ theo tham số ảnh hưởng bước $t=5$ bị nhân với
$\sigma_{\max}(\mathbf{W}_{hh})^{55}$. Với $\sigma_{\max} = 0.9$, hệ số là
$0.9^{55}\approx 3.0\times10^{-3}$; với $\sigma_{\max} = 0.7$, là
$0.7^{55}\approx 3.0\times10^{-9}$.

Điều tinh tế cần nhấn mạnh: mô hình \emph{không phải} không có khả năng biểu diễn
phụ thuộc đó. Trạng thái ẩn hoàn toàn đủ dung lượng mang thông tin ``Huế'' qua
$55$ bước. Vấn đề là \textbf{tín hiệu học} không truyền ngược được quãng đường đó,
nên mô hình không bao giờ \emph{khám phá ra} rằng nó nên mang thông tin ấy. Một
lần nữa, đây là thất bại tối ưu hoá, không phải thất bại biểu diễn --- đúng như
sự phân biệt đặt ra ở Mục~\ref{sec:tongquan}.

\subsection{Băng chuyền sai số hằng: suy diễn thay vì khẳng định}

LSTM thay phép cập nhật \emph{nhân} của~\eqref{eq:rnn} bằng một phép cập nhật
\emph{cộng} trên ô nhớ $\mathbf{c}_t$:
\begin{equation}
  \mathbf{c}_t
  = \mathbf{f}_t \odot \mathbf{c}_{t-1}
  + \mathbf{i}_t \odot \tilde{\mathbf{c}}_t ,
  \label{eq:lstmcell}
\end{equation}
với $\mathbf{f}_t$ là cổng quên, $\mathbf{i}_t$ cổng vào và
$\tilde{\mathbf{c}}_t$ ô nhớ ứng viên.

Lấy Jacobian đầy đủ của~\eqref{eq:lstmcell}. Ba cổng đều là hàm của
$\mathbf{h}_{t-1}$, mà $\mathbf{h}_{t-1} = \mathbf{o}_{t-1}\odot
\tanh(\mathbf{c}_{t-1})$ lại phụ thuộc $\mathbf{c}_{t-1}$. Vậy nên
\begin{equation}
  \frac{\partial\mathbf{c}_t}{\partial\mathbf{c}_{t-1}}
  = \underbrace{\operatorname{diag}\!\left(\mathbf{f}_t\right)}_{\text{đường trực tiếp}}
  \;+\;
  \underbrace{\mathbf{R}_t}_{\text{đường qua các cổng}} ,
  \label{eq:lstmjac}
\end{equation}
trong đó
\[
  \mathbf{R}_t
  = \operatorname{diag}\!\left(\mathbf{c}_{t-1}\right)
    \frac{\partial\mathbf{f}_t}{\partial\mathbf{c}_{t-1}}
  + \operatorname{diag}\!\left(\tilde{\mathbf{c}}_t\right)
    \frac{\partial\mathbf{i}_t}{\partial\mathbf{c}_{t-1}}
  + \operatorname{diag}\!\left(\mathbf{i}_t\right)
    \frac{\partial\tilde{\mathbf{c}}_t}{\partial\mathbf{c}_{t-1}} .
\]

\begin{notebox}
$\mathbf{R}_t$ \textbf{không} bằng không và không được phép lặng lẽ bỏ đi. Nhiều
trình bày phổ thông viết $\partial\mathbf{c}_t/\partial\mathbf{c}_{t-1} =
\operatorname{diag}(\mathbf{f}_t)$ rồi kết luận ngay --- đó là một xấp xỉ, không
phải một đồng nhất thức. Lập luận đúng phải nêu điều kiện để $\mathbf{R}_t$ nhỏ.
\end{notebox}

Mỗi số hạng của $\mathbf{R}_t$ chứa một đạo hàm cổng, mà các cổng đều là sigmoid
hoặc $\tanh$ của một hàm của $\mathbf{c}_{t-1}$ đi qua
$\tanh(\mathbf{c}_{t-1})$. Khi cổng quên bão hoà mở --- $\mathbf{f}_t\to\mathbf{1}$,
tức tiền kích hoạt của nó nằm sâu ở nhánh dương --- thì
$\partial\mathbf{f}_t/\partial\mathbf{c}_{t-1}\to\mathbf{0}$ theo đúng
\eqref{eq:expdecay}; đồng thời nếu $\mathbf{c}_{t-1}$ đã lớn thì
$\tanh'(\mathbf{c}_{t-1})\to0$ làm nhỏ cả hai số hạng còn lại. Trong chế độ đó
$\mathbf{R}_t\approx\mathbf{0}$ và
\begin{equation}
  \prod_{t=2}^{T}\frac{\partial\mathbf{c}_t}{\partial\mathbf{c}_{t-1}}
  \;\approx\;
  \prod_{t=2}^{T}\operatorname{diag}\!\left(\mathbf{f}_t\right)
  \;\approx\;
  \mathbf{I}.
  \label{eq:cec}
\end{equation}
Gradient đi qua $T$ bước thời gian mà \emph{không} chịu suy giảm cấp số nhân. Đây
là \emph{băng chuyền sai số hằng} (constant error carousel).

\subsubsection{Vì sao RNN thuần không thể làm được điều tương tự}

Đối chiếu với tích tương ứng của RNN thuần,
$\prod_{t=2}^{T}\operatorname{diag}(1-\mathbf{h}_{t-1}^2)\mathbf{W}_{hh}^{\top}$,
câu hỏi tự nhiên là: có chọn được tham số nào khiến tích đó bằng $\mathbf{I}$
không? Câu trả lời là không, và lý do mang tính cấu trúc chứ không phải may rủi.

Đặt $\mathbf{W}_{hh} = \mathbf{I}$ --- lựa chọn thuận lợi nhất có thể. Tích trở
thành $\prod_t\operatorname{diag}(1-\mathbf{h}_{t-1}^2)$. Nhưng
$1 - h^2 < 1$ ngặt với mọi $h\neq0$, nên tích này co lại nghiêm ngặt trừ khi
$\mathbf{h}_t = \mathbf{0}$ với mọi $t$ --- tức mạng không mang thông tin nào cả.
Ta gặp một lưỡng nan thực sự:

\begin{guidelinebox}{Lưỡng nan của RNN thuần}
Hoặc trạng thái ẩn khác không --- và khi đó $\operatorname{diag}(1-\mathbf{h}_t^2)$
co chặt, gradient tiêu biến theo cấp số nhân; hoặc trạng thái ẩn bằng không ---
và khi đó mạng không biểu diễn gì. \textbf{Không tồn tại} bộ tham số nào cho phép
RNN thuần vừa mang thông tin vừa truyền gradient không suy giảm qua nhiều bước.
LSTM thoát khỏi lưỡng nan này bằng cách tạo một đường dẫn \emph{cộng tính} đi
vòng qua phi tuyến: trong~\eqref{eq:lstmcell}, $\mathbf{c}_{t-1}$ vào thẳng
$\mathbf{c}_t$ không qua một $\tanh$ nào.
\end{guidelinebox}

Về mặt nguyên lý, \eqref{eq:cec} và kết nối tắt của ResNet là cùng một thủ pháp
đặt ở hai trục khác nhau. ResNet biến Jacobian theo \emph{chiều sâu} thành
$\mathbf{I}+\mathcal{F}'$; LSTM biến Jacobian theo \emph{thời gian} thành
$\operatorname{diag}(\mathbf{f}_t)+\mathbf{R}_t$ với $\mathbf{f}_t\approx\mathbf{1}$.
Trong cả hai trường hợp, một số hạng \emph{bậc không} gần bằng $\mathbf{I}$ được
cài vào tích, và chính số hạng đó gánh gradient đi xa. Xem
Mục~\ref{sec:tongket}.
